\documentclass[output=paper]{langscibook}
\ChapterDOI{10.5281/zenodo.22880298}
\title{\textit{Vagana haopi} – Catching wrasses by hand}
\author{Helen Kobaa Magum and Shalom Magum}
\abstract{}
\begin{document}
\maketitle\label{chap:35}
\begin{figure}[H]
    \captionsetup[subfigure]{oneside,margin={0cm,0cm}}
    \subfloat[]{\includegraphics[trim=0 0 0 0,clip,width=.6\linewidth]{figures/35_1.png}\label{fig:35_1}}\hfil
    \subfloat[]{\includegraphics[trim=0 0 0 0,clip,width=.35\linewidth]{figures/35_2.jpg}\label{fig:35_2}}
    \caption{(a) \textit{Mataberam} -- Rockmover Wrasse, (b) \textit{Haopi Suuna} -- Genuine Wrasse}
\end{figure}
\begin{xltabular}{\textwidth}{lXX}
	(\ref{ex:text28-1})&\textit{O haopi a hena na iana vai to teitei batana paano kasuana.}&\textit{Haopi} (Wrasse) is the name of a fish that stays underground.\\
	(\ref{ex:text28-2})&\textit{Na tei me nana bono buaku o kaku haopi.}&There are two kinds of wrasse.\\
	(\ref{ex:text28-3})&\textit{O vuaba vai o beera na dao raara bono mataberam.}&The big one is called \textit{mataberam} (Rockmover Wrasse).\\
	(\ref{ex:text28-4})&\textit{O vuaba vai o varutaa na dao raara bono haopi suuna.}&The smaller one is called \textit{haopi} \textit{suuna} (genuine wrasse).\\
	(\ref{ex:text28-5})&\textit{Evehee o hena beera o haopi.}&But the general name is wrasse.\\
	(\ref{ex:text28-6})&\textit{Beara tau vagana haopi, eara repaa tara rakerake bata o sii vahara vasu beori ravun vaa tea vaakororovi rori.}&When we are about to catch wrasses, we search for little stones being heaped up in a circle.\\
	(\ref{ex:text28-7})&\textit{Beara tara ahe eve, eara toro huu hopo ni nao a sinaka,}&When we spot it, we must push the fishing-stick into it,\\
	(\ref{ex:text28-8})&\textit{ei pasi pahin tahi ahe koa vuara.}&this (means) we will immediately spear it.\\
	(\ref{ex:text28-9})&\textit{Eara repaa kasi, are kasi antee beara kasi ahe komana e.}&Then we will dig and dig until we recognize it.\\
	(\ref{ex:text28-10})&\textit{Ato nie teo buo ketaa nae, eara repaa rahi tavusu e.}&(We) hold it on its two gills, and then we pull it out.\\
	(\ref{ex:text28-11})&\textit{Ei asun vamate vai ara.}&Now we kill it.\\
	(\ref{ex:text28-12})&\textit{Evehee na teitei batana a sii vapara nae.}&But usually there is also its partner.\\
	(\ref{ex:text28-13})&\textit{Goe tea nao vananaoana vira bean asun vakavara o peho.}&Don't go too far away when you have killed one.\\
	(\ref{ex:text28-14})&\textit{Tara rakerake vaha koa o meho rihinae.}&Just search for the other one nearby.\\
	(\ref{ex:text28-15})&\textit{Ean paa kasi e.}&You dig it out.\\
	(\ref{ex:text28-16})&\textit{Eara samin doana haara, na pahin kasi ona bata koa raara.}&We never fail, we immediately dig (and) catch (it).\\
	(\ref{ex:text28-17})&\textit{Eara toro matatopo bata koara to mene rosin tavusu maa.}&We must be prepared that it might come out (and) escape.\\
	(\ref{ex:text28-18})&\textit{Kikira vakikis bata koa ni nom eve.}&Just keep carefully watching it.\\
	(\ref{ex:text28-19})&\textit{Eara mee toro goe tea ato o haopi en vihi nae, ei kou to mene dohi ore paa rosin.}&Also, we must not hold the wrasse on its tail because it might slip (out of your hands) and escape. \\
	(\ref{ex:text28-20})&\textit{Evei o sii vahutate vaa tea vagana haopi.}&This is the little description of catching wrasses.\\
\end{xltabular}
\end{document}